\documentclass[english,twoside, openright, hidelinks, 11 pt, class=report,crop=false]{standalone}
\input{../bm}
\input{../../preamb}


\begin{document}

\subsection*{Kapittel \ref{Alg}}
\opr{opgalgfullkvad}
\abch{
	\item $ (x+3)^2 $
	\item $ (b+7)^2 $
	\item $ (a-1)^2 $
	\item $ \left(k-\frac{1}{3}\right)^2 $
	\item $ \left(c-\frac{1}{4}\right)^2 $
	\,\item $ \left(y+\frac{3}{7}\right)^2 $
}

\opr{opgalgfullkvad2}
\abch{
	\item $ (5a+9)^2 $
	\item $ (3b+2)^2 $
	\item $ (8c-1)^2 $
	\item $ \left(\frac{1}{2}d+\frac{3}{4}\right)^2 $
	\item $ \left(\frac{1}{5}e+\frac{2}{7}\right)^2 $
	\item $ \left(\frac{9}{8}-\frac{5}{3}\right)^2 $
}

\opr{opgalgkongsetn}
\abch{
	\item \selos
	\item \selos
}

\opr{1tv22d1opg}
$s=5$, $r=25$


\opr{1tv21d1opg8}
$s=2$, $r=64$, $t=8$
\[ 4x^2+16x+r=(sx+t)^2 \]

\opr{1tv24d1opg3}
$(a+b)(a-b)=a^2-b^2$


\opr{opgalgbrokforenkl} 
\abch{
	\item \selos
	\item $\frac{5}{13}(4+\sqrt{3})$
	\item \selos
}

\opr{opgalgledd}
\abch{
	\item $ (x+3)^2-16 $
	\item $ (x-4)^2-36 $
	\item $ (x+6)^2+81 $
	\item $ (2x-3)^2+2 $
	\item $ (7x+9)^2-1$
}

\opr{algopga1a2}
Så lenge vi søker heltall, er det bare noen få heltall som vil oppfylle kravet $ a_1a_2=c $, mens uendelig mange heltall oppfyller kravet $ a_1+a_2=b $.

\opr{opgalgfakt}
$ (x+3)^2-16=(x+7)(x-1) $, $ (x-4)^2-36=(x+2)(x-10) $

\opr{opgalgfullkvad3}
\abch{
	\item $ (x-5k)^2$
	\item $ (y+4z)^2 $
	\item $ (a-10q)^2 $
}

\opr{1TV23D1opg3}
$a=2$ og $b=6$ eller $a=-2$ og $b=-6$

\opr{1TH21 og 1TV25}{opgalgeksuneq}
\abch{
	\item $-4<x<2$
	\item $-2<x<6-2 < x < 6-2<x<6$
}

\opr{opgalgulikgraf}
\abc{
	\item $x\in (-\infty, 3) \cup (7, \infty)$
	\item Løs ulikheten ved hjelp av faktorisering.
}

\opr{1TH21D1opg3}
\abch{
	\item $\frac{2(x+1)}{x-1}$
	\item $\frac{2x}{x+3}$
}


\opr{opgalgulikbrok}
\abch{
	\item Man kan ikke vite om fellesnevneren er positiv eller negativ, og dermed vet man ikke om ulikhetstegnet må snus eller ikke.
} \\
\abchs{2}{
\item $x\in (-5,-4)\cup (-2,\infty)$
}

\opr{opgalgligukonst} 
\selos

\opr{opgalgligukonst2}
\abch{
	\item $ x=0 \vee x=2$
	\item $ x=0 \vee x=-9 $ 
	\item $ x=0 \vee -\frac{2}{7} $
	\item $ x=0 \vee x=\frac{8}{9} $
}

\opr{1TV23D1opg2}
$x=-2$ og $x=4$

\opr{opgalgabc}
\abch{
	\item $ x=3\vee x=-5 $
	\item $ x=7\vee x=-10 $ 
	\item $ x=\frac{-5\pm\sqrt{53}}{2} $
	\item $ x=\frac{1\pm\sqrt{5}}{2} $
	\item $ x=1\pm\sqrt{10} $	
	\item $ x=2\pm 2\sqrt{2} $ 	
	\item $ x=1\vee x=-\frac{7}{5} $
	\item $x=\pm\frac{\sqrt{6}}{2} $
	\item $x=2\pm\frac{\sqrt{33}}{3} $ 
}



\opr{1TH21D1opg4}
$f(x)=-4x^2+4x+12$

\opr{1tv21opg4}
$k=\pm\frac{3}{2}$


\opr{opgalgsym}
$x=-1$

\opr{1TH21D1opg5} 
\abc{
	\item \selos
	\item \selos
}

\opr{opgalgabcomv}
\abch{
	\item 0 for begge. $t$ og $s$ er løsningene av likningen $f(x)=0$.
	\item \selos
	\item $\frac{\sqrt{b^2-4ac}}{a}$. Dette er distansen mellom nullpunktene.
	\item $\frac{-b}{2a}$. Dette er $x$-verdien til midtpunktet mellom nullpunktene.
	\item $\frac{c}{a}$. $a=5$ og $c=15$
}


\opr{algopgpoldiv}
Utfør polynomdivisjon på uttrykkene \os
\abch{
	\item $ x-2 $
	\item $ x^2-1 $
	\item $ x-3 $
	\item $x+4+\frac{9}{x-2}$
}

\opr{1tv24d1opg2}
\selos

\opr{1tv22d1opg6}
$f(3)=0$, da er $f$ delelig med $x-3$.

\opr{algopgpolfakt}
\abch{
	\item $ P=(x-3)(x-4)(x+7) $
	\item $ P=(x+1)(x+3)(x+5) $
	\item $ P=(x+1)(x+6)\left(x+\frac{7}{2}\right) $
}

\opr{1tv21d1opg3}
$k=-2$

\opr{opgalgfind}
$(-4, 0)$ og $\left(\frac{4}{3}, \frac{32}{3}\right)$


\opr{R1H23D1opg2}
$2\ln e<3\log_{10} 70<e^{3\ln 2}$

\opr{r1h22d1opg3}
$3\sqrt{11}$ og $4\ln 9$

\opr{opgalgpot1}
Løs likningen. \os
\abch{
	\item $ x=\frac{\ln 2}{\ln 5} $
	\item $ x=\frac{\ln 9}{\ln 8} $
	\item $ x=\frac{1}{\ln 2}\ln\frac{19}{11} $
}


\opr{opgalgpot2}
\selos

\opr{opgzeros}
\abch{
	\item $x=\frac{1}{2} $
	\item $x=-3$ og $x=1$
}

\opr{opgalgsub}
\abch{
	\item $ x=e^2\vee x=e^3$
	\item $x=e^{-7}\vee x=e^{10}$
	\item $ x=\ln 3 $ 
	\item $ x=\ln 2$
	\item $x=\ln 2$
	\item $x=\frac{1}{100}\vee x=10000$
	\item $x=\log 3$ 
	\item $x=e^{-2}\vee e^3$
}

\opr{r1h25d1opg2}
$a=4$

\opr{opgalglog}
\abch{
\item $ \frac{1}{2}(6+e^2) $
\item $ x=2$
\item $ x=\frac{3}{2} $ 
\item $x=e^{-x} $
}

\grubr{opgalglogxlnx}
\selos

\grubr{r1h24d2opg4}
$a=5$

\grubr{t1h23d1opg2}
$x=-3$, $x=-1$ og $x=2$


\grubr{1tv24d1opg5}
\abch{
	\item $f(x)=-2(x+3)(x-4)$.
	\item $-2<x<3$. 
}

\grubr{1tv25d1opg3}
\selos

\grubr{opgalgsqrt27}
$x=3$ og $y=12$

\grubr{opgalgfaktfourterms}
\abch{
\item $ (x+5y)(x-4y) $
\item $ (a-2b)(a-7b)$
\item $ (k^2-y)(9k^5-y)$
}

\grubr{opgalgtre}
$x=-2$, $x=-10$ og $x=1$

\grubr{opgalgkvadsqrt}
\abch{
	\item $ (\sqrt{7}-\sqrt{3})^2 $
	\item $ (\sqrt{11}+\sqrt{2})^2 $
	\item $ (\sqrt{6}+\sqrt{2})^2 $
	\item $ (\sqrt{35}-\sqrt{7})^2 $
}

\grubr{opgfunkfindx}
$c=6$

\grubr{opgalgq3minp3}
\abch{
	\item $ (q-p)(q^2+pq+p^2) $
	\item \selos
}

\grubr{opggeoher}
\selos

\grubr{opgaldrdh}
$(r^2 + 2b d)(2b d - 2d r-r^2)$

\grubr{opggeokvadsym}
\selos

\grubr{opgalgsym2}
\selos

\grubr{1th24d1opg2}
$(-4, -1)$

\grubr{1th24d1opg3}
$-2<x<1$ og $x<-6$\;\;  \includegraphics[scale=0.25]{\figp{1th24d1opg3_lf}}


\grubr{1th21d1opg3}
$x=-4\vee x=1$

\grubr{opgbevisabc}
\selos




\end{document}